\section{Threats to Validity}
\label{sec:threats}

\paragraph{Construct validity}
Output coverage $\OCov_s$ measures adequacy on the abstracted output space, so its
meaningfulness depends on the partition scheme $\{\pi_j\}$. A poorly chosen
abstraction could inflate or deflate coverage. We mitigate this by defining
partitions from established behavioural properties (calibration, fairness,
monotonicity, robustness, measurement fidelity) rather than ad hoc bins, and by
showing (RQ4) that conclusions are stable as granularity varies. Fault detection
is measured against seeded behavioural faults expressed as output-interaction
signatures; we seed rare-but-reachable signatures to avoid trivially detectable
faults, and we report the tie with DeepCT rather than selecting a fault set that
favours our method.

\paragraph{Internal validity}
Randomness enters through model training, feasibility probing, and the inverse-map
optimisers. We control it with fixed seeds and report 30-run distributions with
non-parametric tests (Mann--Whitney $U$, bootstrap CIs, Cliff's $\delta$) rather
than single-run point estimates. The prioritised suite is truncated at the
coverage plateau, a deterministic rule, so suite-size comparisons across methods
use matched budgets.

\paragraph{External validity}
Our study covers four domains but a limited number of systems and datasets per
domain. The quantum SUT is a noisy statevector simulation rather than hardware,
and the vision/NLP inverse maps operate in latent (PCA / SVD) spaces with
decode-to-manifold steps, which bound the realism of synthesised inputs. We do not
claim that $\OCov_s$ correlates with real-world defect discovery beyond the seeded
faults; establishing that link on production systems is future work. The results
should be read as evidence that output-interaction coverage is \emph{attainable}
and \emph{discriminating} across diverse system types, not as a claim of universal
superiority.

\paragraph{Conclusion validity}
We use non-parametric tests that make no normality assumption and report effect
sizes alongside $p$-values. Where \rnwise{} does not dominate---fault detection
vs.\ DeepCT ($p=0.92$) and tuple efficiency vs.\ DeepCT---we say so explicitly.
The headline claim is confined to what the data support: unique, statistically
significant \emph{output coverage}, at low search cost, across domains.
